\subsection{Why uniform verified replay prevents extinction}
\label{app:theory-replay}

A persistent replay signal becomes a retention guarantee when vanishing
probability would force a bounded loss to diverge. The next argument makes
that loss barrier explicit, then shows what changes with unequal weights
and incomplete coverage.

Let $\mathcal B\subseteq\mathcal C$ be a nonempty fixed bank of $k$ modes that
the policy has produced and the validator has accepted. Uniform verified replay
uses one exemplar per banked mode and, in the categorical reduction, minimizes
\begin{equation}
  R_{\mathcal B}(z)
  =-\frac1k\sum_{b\in\mathcal B}\log p_b.
  \label{eq:verified-replay-loss}
\end{equation}
This is ordinary cross-entropy against the distribution that assigns mass
$1/k$ to each banked mode and zero elsewhere. It is not conditioned on the
bank: decreasing~\eqref{eq:verified-replay-loss} can raise total bank mass as
well as redistribute mass within the bank.

In an averaged continuous-time model, absorb a fixed round-robin replay
frequency and loss scale into an effective dose $\rho>0$. This assumes a
fixed eligible-bank set and constant relative weighting; it does not identify
finite alternating updates with simultaneous gradient flow. The prompt-isolated categorical mean flow for any of the three fresh
binary-reward objectives above plus verified replay is
\begin{equation}
  \dot z
  =c_G(P)\nabla_zP-\rho\nabla_zR_{\mathcal B}(z).
  \label{eq:replay-mean-flow}
\end{equation}
Let $\Psi(P)=\int_0^P c_G(s)\,ds$. For the finite-group estimators in Lemmas~\ref{lem:grpo-mean}
and~\ref{lem:maxrl-mean}, $c_G$ has a finite positive extension to $[0,1]$;
in particular, Dr.GRPO has $c_G=(G-1)/G$ and MaxRL has
$c_G=[1-(1-P)^{G-1}]/P$. Define
\begin{equation}
  F_{\mathcal B}(z)=\rho R_{\mathcal B}(z)-\Psi(P(z)).
  \label{eq:replay-potential}
\end{equation}
Then~\eqref{eq:replay-mean-flow} is exactly gradient flow
$\dot z=-\nabla_zF_{\mathcal B}$.

\begin{theorem}[Verified replay prevents banked-mode extinction]
\label{thm:replay-retention}
Suppose that after time $T$ the nonempty verified bank $\mathcal B$ is fixed,
fits within replay capacity, and receives a fixed positive effective replay
dose $\rho$. Under~\eqref{eq:replay-mean-flow}, define
\[
  C_T
  =R_{\mathcal B}(z(T))
   +\frac{\Psi(1)-\Psi(P(T))}{\rho}.
\]
Then for every $b\in\mathcal B$ and every $t\ge T$,
\begin{equation}
  p_b(t)\ge \exp(-kC_T)>0.
  \label{eq:replay-retention-bound}
\end{equation}
Thus each banked probability is uniformly bounded away from zero for all
$t\ge T$, including asymptotically, even when the bank does not cover every
correct mode. This is stronger than finite-time positivity, which already
follows from finite softmax logits without replay.
\end{theorem}

\begin{proof}
Along the flow,
\[
  \frac{d}{dt}F_{\mathcal B}(z(t))
  =-\lVert\nabla_zF_{\mathcal B}(z(t))\rVert_2^2\le0.
\]
Since $P(t)\le1$ and $\Psi$ is increasing,
\[
  R_{\mathcal B}(z(t))
  \le R_{\mathcal B}(z(T))
     +\frac{\Psi(1)-\Psi(P(T))}{\rho}
  =C_T.
\]
Every $-\log p_b$ is nonnegative and their sum is
$kR_{\mathcal B}$. Therefore
$-\log p_b(t)\le kC_T$ for every banked mode, which is exactly
Equation~\eqref{eq:replay-retention-bound} for every time after $T$.
\end{proof}

\paragraph{Worked example: a rare mode cannot hide inside a bounded loss.}
In \hyperref[ex:theory-running]{the running example}, bank the first and third correct modes,
whose probabilities are $.30$ and $.05$. Their uniform replay loss is
$R_{\mathcal B}=-\tfrac12\log(.30\cdot.05)\approx2.09985$.
For the illustrative Dr.GRPO choices $G=4$ and $\rho=.30$,
$C_T=R_{\mathcal B}+(.75)(1-.50)/.30\approx3.34985$, giving a floor
$p_b\ge e^{-2C_T}\approx0.00123127$. The point is the constraint on the
loss: if either banked probability fell to $10^{-4}$, its term alone would
contribute $-\tfrac12\log(10^{-4})\approx4.60517>C_T$.
The other term is nonnegative and cannot cancel this increase. These are
hypothetical categorical values; the argument explains the barrier without
making its conservative floor an experimental prediction.

This is a qualitative positivity statement, not a useful numerical guarantee.
If a prompt receives fresh updates with frequency $\nu_x>0$ and replay with
frequency $\mu_x>0$, normalizing time by its fresh-update frequency gives
$\rho=\rho_{\rm raw}\mu_x/\nu_x$. Equal fresh and replay frequencies cancel;
round-robin scheduling alone does not imply dilution of the relative dose.
As an illustrative equal-frequency categorical calculation, the registered
$G=16$ coefficient and per-rollout scale give
$\rho=.10(15/256)=5.86\times10^{-3}$. With $k=16$, the worst-case Dr.GRPO
potential term alone permits a bound on the order of $\exp(-2560)$ (and
$R_{\mathcal B}(z(T))$ can only weaken it). The bound is therefore vacuous at
machine and experimental scales; it establishes only that an idealized
continuous-time, fixed-bank coordinate cannot vanish asymptotically.

\paragraph{The retention bound also holds under smooth parameterization.}
The same argument applies to $\dot\theta=-\nabla_\theta F_{\mathcal B}(z(\theta))$:
\[
 \frac{d}{dt}F_{\mathcal B}(z(\theta(t)))
 =-\lVert J_\theta^\top\nabla_zF_{\mathcal B}\rVert_2^2\le0.
\]
Thus Equation~\eqref{eq:replay-retention-bound} holds regardless of parameter
count for this fixed-bank, prompt-isolated parameter gradient flow. It does
not require the shared-correctness geometry in
Theorem~\ref{thm:shared-correctness}. This extends the qualitative retention
bound, not convergence to a uniform distribution: a restricted or degenerate
parameterization can stall. The next corollary retains independent logits.

\begin{corollary}[Conditional full-coverage categorical limit]
\label{cor:replay-no-collapse}
Assume the conditions of Theorem~\ref{thm:replay-retention} and that by time
$T$ the bank covers the full correct support, $\mathcal B=\mathcal C$. Let
$u_c=1/m$ on $\mathcal C$. Then the categorical mean flow converges to
\[
  P(t)\longrightarrow1,
  \qquad q(t)\longrightarrow u.
\]
Consequently
\[
  \min_{c\in\mathcal C}q_c(t)\longrightarrow\frac1m,
  \qquad
  \mathcal H(q(t))\longrightarrow\log m,
\]
and conditional correct-mode collapse is impossible.
\end{corollary}

\begin{proof}
When $\mathcal B=\mathcal C$,
\begin{equation}
  R_{\mathcal C}
  =\log m-\log P+\mathrm{KL}(u\|q).
  \label{eq:replay-mass-uniformity-decomposition}
\end{equation}
Thus one uniform verified-likelihood loss already contains both a verified-mass
term $-\log P$ and a within-support uniformity term $\mathrm{KL}(u\|q)$. The
potential~\eqref{eq:replay-potential} becomes
\[
  F_{\mathcal C}
  =\rho\log m+\rho\mathrm{KL}(u\|q)
   -\rho\log P-\Psi(P).
\]
For fixed $P$, its unique minimizer is $q=u$. Its derivative in $P$ is
$-\rho/P-c_G(P)<0$, so its unique minimizing distribution on the closed
simplex is $P=1,q=u$,
placing all limiting probability uniformly on the correct support and zero
mass on every incorrect category in the categorical limit.

For completeness, Theorem~\ref{thm:replay-retention} keeps every correct
coordinate away from zero. The decreasing potential gives
$\int_T^\infty\lVert\nabla F_{\mathcal C}\rVert_2^2dt<\infty$.
The finite-group $c_G$ is smooth on $[0,1]$; softmax derivatives and the
cross-entropy gradient are bounded. Thus the Hessian of $F_{\mathcal C}$ and
the velocity $\dot z=-\nabla F_{\mathcal C}$ are bounded, making
$\nabla F_{\mathcal C}(z(t))$ uniformly continuous. Together with square
integrability, this implies $\nabla F_{\mathcal C}\to0$: recurring gradients
of fixed positive size would persist over intervals of a uniform minimum
length, contradicting the finite energy integral. At any limiting
distribution the correct-coordinate stationarity equations are
\[
  0=\rho(p_c-1/m)-c_G(P)p_c(1-P),
\]
so all $p_c$ are equal. Substituting $p_c=P/m$ gives
$0=-(1-P)(\rho+c_G(P)P)/m$, hence $P=1$. Each incorrect coordinate satisfies
$0=p_a(\rho+c_G(P)P)$ and therefore has $p_a=0$. The limiting distribution is
unique, proving convergence from every initialization covered by the corollary.
\end{proof}

\paragraph{Positive replay weights and exemplar length.}
Uniformity of the target is separate from protection against extinction.
For fixed weights $w_b>0$ with $\sum_{b\in\mathcal B}w_b=1$, define
$R_w=-\sum_b w_b\log p_b$ and use replay coefficient $\rho_w>0$.
The same potential argument gives, for all $t\ge T$,
\[
 p_b(t)\ge\exp(-C_{T,w}/w_b),\qquad
 C_{T,w}=R_w(z(T))+\frac{\Psi(1)-\Psi(P(T))}{\rho_w}.
\]
If $\mathcal B=\mathcal C$, then
$R_w=H(w)-\log P+\mathrm{KL}(w\|q)$, and the independent-logit limit is
$P\to1$, $q\to w$. The preceding stationarity proof applies with $1/m$
replaced by $w_c$: summing the correct-coordinate equations forces $P=1$,
then each equation forces $p_c=w_c$.

This distinction matters for the implemented mean-token exemplar loss.
Even if each key were represented by exactly one categorical exemplar of
fixed length $L_b$, its loss would be $A R_w$, where
\[
 A=\frac1k\sum_{b\in\mathcal B}\frac1{L_b},\qquad
 w_b=\frac{L_b^{-1}}{\sum_{j\in\mathcal B}L_j^{-1}}.
\]
Its effective categorical replay coefficient would be $\rho_w=\rho_{\rm LM}A$.
The full-coverage limit would therefore weight shorter exemplars more, and
would be uniform only for equal lengths. For an actual language model,
$\pi_\theta(e_b\mid x)$ is only one contribution to the total key probability.
Equal averaging of per-token exemplar scores does not imply uniform learned
key probabilities; the categorical uniformity corollary is not that claim.

\paragraph{Worked example: equal key sampling need not imply a uniform target.}
For three singleton categorical exemplars of lengths $10,20,40$, equal
averaging of mean-token losses gives $A=7/120$ and
$w=(4/7,2/7,1/7)$. All three target weights are positive, so each mode is
protected, but the full-coverage categorical limit assigns the shortest
exemplar four times the probability of the longest. Equal lengths would
instead give $(1/3,1/3,1/3)$. This example concerns the loss's induced
categorical target, not an asserted distribution over actual LM keys.

The premise $\mathcal B=\mathcal C$ is not a description of the released
Python-factors experiments. Their prompts have 16--3,600 valid execution modes
(mean 229.44), while the replay bank has capacity 16. Thus the corollary is
inapplicable to Python factors as a domain-wide guarantee. More generally it
applies only to prompts whose entire correct support fits in the bank and has
actually been discovered; neither condition follows from replay.

\begin{lemma}[Fixed-bank limits and deterministic finite steps]
\label{lem:replay-finite-steps}
Retain finite independent Euclidean logits, a nonempty fixed bank
$\mathcal B\subseteq\mathcal C$, fixed positive normalized weights $w_b$,
and $\rho_w>0$. Let $F=\rho_wR_w-\Psi(P)$, with
$\Psi\in C^2([0,1])$ and $\Psi'=c_G\ge0$. Set
\[
 M_1=\sup_{[0,1]}|\Psi'|,\qquad
 M_2=\sup_{[0,1]}|\Psi''|,\qquad
 L=\frac{\rho_w}{2}+\frac{M_1}{2}+\frac{M_2}{8}.
\]
For constant $0<\eta<2/L$, deterministic updates
$z_{n+1}=z_n-\eta\nabla F(z_n)$ preserve, for every $n\ge0$,
\[
 p_b(z_n)\ge \exp(-C_{0,w}/w_b),\qquad
 C_{0,w}=R_w(z_0)+\frac{\Psi(1)-\Psi(P(z_0))}{\rho_w}.
\]
Moreover, $p(z_n)\to\bar w$, where $\bar w_b=w_b$ on the bank and
$\bar w_a=0$ elsewhere. The same probability limit holds for the corresponding
continuous mean flow. The logits need not converge: zero limiting probabilities require divergent logit differences relative to banked coordinates.
\end{lemma}

\begin{proof}
For a logit direction $u$, let $U=u_a$ and $Y=\mathbf1[a\in\mathcal C]$
under $a\sim p$. Differentiating along that direction gives
\[
 u^\top\nabla^2P\,u
 =\mathbb E[(Y-P)(U-\mathbb EU)^2],\qquad
 u^\top\nabla^2R_w\,u=\operatorname{Var}(U).
\]
Since $|Y-P|\le1$, the first quadratic form is bounded in absolute
value by $\operatorname{Var}(U)$. The bounded-range variance inequality gives
\[
 \operatorname{Var}(U)\le\frac{(\max_a u_a-\min_a u_a)^2}{4}
 \le\frac{\lVert u\rVert_2^2}{2}.
\]
Both Hessians therefore have operator norm at most $1/2$.
Also $\lVert\nabla P\rVert_2^2\le2P^2(1-P)^2\le1/8$.
The chain rule therefore bounds $\lVert\nabla^2F\rVert_{\rm op}\le L$
globally. The descent lemma yields
\[
 F(z_{n+1})\le F(z_n)
 -\eta(1-L\eta/2)\lVert\nabla F(z_n)\rVert_2^2.
\]
Energy monotonicity gives the preceding replay-loss bound and hence the
stated retention floor. Since $F\ge-\Psi(1)$, the squared gradients are
summable and $\nabla F(z_n)\to0$. Its coordinate expression,
\[
 [\nabla F]_a=\rho_w(p_a-\bar w_a)
 -c_G(P)p_a(\mathbf1[a\in\mathcal C]-P),
\]
extends continuously to the closed probability simplex. Every subsequential
probability limit is therefore stationary. On an incorrect coordinate,
stationarity gives $(\rho_w+c_G(P)P)p_a=0$, forcing all incorrect mass to
zero and $P=1$. The correct-coordinate equations then give $p_a=\bar w_a$.
Compactness and this unique limit prove convergence of the whole probability
sequence. For the mean flow, bounded Hessian and velocity together with
energy descent give vanishing gradient as in the preceding corollary;
the same stationarity argument applies.
\end{proof}

For Dr.GRPO, $M_1=(G-1)/G$ and $M_2=0$. For the implemented finite-group
MaxRL estimator, $M_1=G-1$ and $M_2=(G-1)(G-2)/2$, so one sufficient smoothness
constant is $L=\rho_w/2+(G-1)(G+6)/16$. These bounds concern categorical logit
steps, not the neural optimizer's learning rate. With an incomplete frozen
bank, the categorical limit assigns zero mass even to correct modes outside
the bank: fixed-bank retention is not a guarantee of discovering or preserving
all correct modes. The fixed-bank limit also gives
\[
 \lim_{t\to\infty}\texttt{distinct@}K
 =\sum_{b\in\mathcal B}[1-(1-w_b)^K]\le k,
\]
with the same expression along deterministic iterates. Thus the bank itself
sets an asymptotic support ceiling in this model. The later stochastic and changing-bank results give conditional
retention bounds under additional energy assumptions; they do not extend this
fixed-bank convergence limit to arbitrary neural training.

\paragraph{Worked example: a safe logit step and an incomplete bank.}
Return to the two-mode bank above with $\rho_w=.30$, $G=4$, and Dr.GRPO.
Here $L=.525$, so $\eta=1$ satisfies $0<\eta<2/L\approx3.81$.
One deterministic logit-gradient step changes the running distribution
approximately as
\[
 (.30,.15,.05,.50)\ \longmapsto\ (.3861,.1643,.0632,.3864),
\]
while the potential falls from about $.25496$ to $.09689$.
Both banked modes gain probability. The unbanked second correct mode also
gains on this first step, but the lemma's eventual limit is
$(.50,0,.50,0)$: early improvement of an unbanked mode does not make it
protected. The step size here belongs to the independent-logit example;
it is not a recommendation for the neural optimizer.

